\section{Introduction}
Fine-grained classification is a challenging yet important task that involves
differentiating subtle distinctions in a label set. For instance, in image
classification, classifying cars, planes, and other vehicles is
\textit{coarse-grained} classification, whereas distinguishing
models from cars is \textit{fine-grained} classification.
In NLP, sentiment analysis, which attempts to classify positive/negative
sentiments, is an example of coarse-grained text classification. Human emotions,
however, exhibit more complexity. For example, the six type of basic
emotion~\cite{ekman_basic_emotion} include \textit{happiness, sadness, fear,
disgust, anger}, and \textit{surprise}, and show finer distinctions of positive and
negative classes. Moreover, complex interactions exists in human emotion as
different type of emotions can have subtle differences, for instance
\textit{ashamed} and \textit{guilty}. This makes fine-grained emotion classification (FEC) 
challenging not only
because of the increased number of classes, but also because of the increased similarity between classes.
For instance, the current finest emotion classification datasets
contain up to 27 and 32 classes of
emotion~\cite{rashkin-etal-2019-towards,demszky-etal-2020-goemotions},
respectively. Classes such as \textit{furious} and \textit{angry} in these
fine-grained datasets are far more difficult to differentiate than
\textit{happy} and \textit{sad}. 
\begin{figure}[t!]
\includegraphics[width=\columnwidth]{figures/intro_poincare.png}
\caption{\label{fig:intro_poincare}Properties of hyperbolic space.
The hyperbolic distance between two points grows dramatically as they come close to
the border, converse to the Euclidean distance. Moreover, the hyperbolic
distance is an analog of   % AMH: check
tree distance, which is $d(B,C) \approx d(B,A) + d(A,C)$. }
\end{figure}
However, detecting fine-grained emotion is useful in a variety of situations.
For example, in a dialog generation system, understanding the user's
fine-grained emotion could facilitate more empathy in the responses, which might
increase conversation engagement~\cite{roller-etal-2021-recipes}. 
Emotions also play an important role in people's lives, affecting how they make decisions and how they interact with machines. Therefore, the finer we can classify, the more information we can collect to model users' behaviors.

Existing text classification methods often use pre-trained language models such as
BERT~\cite{devlin-etal-2019-bert} to generate a contextualized representation,
and fine-tune them primarily in Euclidean space for downstream tasks. 
%They are effective because they make use of a large amount of unlabeled data and the self-attention mechanism.
However, such a space is limited, as some confusing pairs have
nearly identical meanings (e.g., \textit{furious} and \textit{angry}), 
%which is hard to tell the difference even by humans \lwku{I feel redundant here}, 
and forcing them to be separated in the latent space may harm the performance
by overfitting on training data.
The complex nature of emotions can be expressed in a hierarchical way, consisting of three levels~\cite{parrott2001emotions}.
For instance, \emph{Joy}, \emph{Contentment}, and \emph{Pleasure} are primary, secondary, and tertiary emotions, respectively.
Meanwhile, learning embeddings in hyperbolic space is becoming more popular due to its superior
ability to capture hierarchical information~\cite{NIPS2017_poincare_hierarchical, 
%poincare_embedding_NIPS2017,  % duplicate ref
ganea2018hyperbolic,chami2019hyperbolic_gcn,liu2019hyperbolic_gnn}.
Figure~\ref{fig:intro_poincare} demonstrates a tree embedded to hyperbolic
space. Tree nodes closer to the root are embedded near the origin,
and nodes closer to leaves are placed closer to the boundary. The
main merit of this space is that as the distance from the origin rises, the
amount of space in hyperbolic space grows
exponentially~\cite{cho2019large,lopez-strube-2020-fully,2021_hyper_survey}.
Intuitively, tree-like structures also expand the number of nodes as the
distance increases from the root, which is consistent with the mathematical basis
of hyperbolic geometry. This is also reflected in 
the hyperbolic distance (Eq.~\ref{eq:poincare_distance}), % AMH: check
which resembles
the distance between two nodes in a tree.

In this work, we propose \modelname, which integrates label embedding trained
in hyperbolic space with a RoBERTa model fine-tuned in Euclidean space.
Specifically, we first learn hierarchical-aware label embeddings in 
hyperbolic space, and then project the representation output by RoBERTa onto
the same space to derive the distance between a text representation and its
corresponding label. This distance is then used to weight standard cross
entropy loss, making the projection of text representation as close to its
label as possible, in hyperbolic space. 
%\lwku{I feel here too detail...should be more conceptual} 
%We discovered that jointly learning representation on both hyperbolic and euclidean space greatly improves the FEC task. 
Results on two challenging datasets,
GoEmotions~\cite{demszky-etal-2020-goemotions} and
EmpatheticDialogs~\cite{rashkin-etal-2019-towards}, demonstrate the superiority of the
proposed model. Also, we find that \wsmodelname performs best when the
label structure is more complex, or the inter-class relationship is more
ambiguous.
To sum up, the contributions of this paper are threefold:
\begin{itemize}
	 \item We leverage the merits of hyperbolic geometry to learn better
	 representations in both hyperbolic and Euclidean space. 
	 \item We propose the novel \wsmodelname framework along with a simple yet
	 effective objective function to address the FEC task.
	 \item Empirically, the proposed model outperforms existing methods, and
	 is even comparable with systems that utilize external knowledge or data
	 augmentation techniques. 
\end{itemize}
